\section{Completed Derivation Graph and Global Closure Theorem}
\label{sec:derivation-graph-verdict}

The PR-I--PR-III construction is a dependency-ordered directed acyclic graph. Its public sector maps and finite-tier rules are inherited relations; its equation-selection coordinates are exactly those listed in Proposition~\ref{prop:pr35-complete-certified-class}.  Every node in the claimed theorem domain is therefore primitive, derived, a proved canonical closure, a fixed inherited map, or an equivalent presentation.  Diagnostic references have no generative edges.  No unresolved selector is an element of the completed admissible class.

\begin{table}[H]
\centering
\caption{Final Uniqueness Classification in the Completed PR Admissible Class.}
\label{tab:pr35-status}
\begin{tabularx}{\textwidth}{@{}p{0.28\textwidth}X@{}}
\toprule
Block & Final result \\
\midrule
PR-I radial coefficients & Sole minimal even-quartic normal form $(a_2,a_4)=(1,3/4)$. \\
PR-I compact projector & Sole minimal required-channel range $\operatorname{span}\{\phi_1,\phi_3\}$. \\
PR-I terminal boundary & Sole signed zero-incidence value $\sigma_8=-1$. \\
PR-II hypercharge & Sole nonzero-charge orbit modulo normalization and up/down relabeling. \\
PR-II photon & Sole null line of the fixed representation-derived outer-product neutral mass operator. \\
PR-II generations & Sole normalized complete/nonredundant incidence orbit $(N_{\rm gen},[\Phi_3])=(3,[I_3])$. \\
PR-III finite-tier maps & $D_1,D_2,C_1,C_2$ and the declared neutrino/strong maps are fixed inherited finite-tier relations, not free selector coordinates. \\
PR-III neutral response & Sole normalized minimal-response tensor $D_3=2\rho P_Z$. \\
Provenance & No diagnostic target is a generative ancestor and no fitted coefficient occurs. \\
Complete admissible quotient & One dependency-ordered equation architecture modulo ledger-preserving equivalence. \\
\bottomrule
\end{tabularx}
\end{table}

\begin{theorem}[Singleton Theorem for the Completed PR Admissible Class]
\label{thm:pr35-completed-selector-singleton}
Let $\mathfrak A_{\rm PR}^{\rm can}$ be the complete class in Proposition~\ref{prop:pr35-complete-certified-class}, and let $\sim$ identify gauge, normalized-basis, coordinate, phase, algebraic, field-redefinition, and unit presentations that preserve the complete public ledger and registered
observable algebra.  Then
\begin{equation}
 \boxed{
 \mathfrak A_{\rm PR}^{\rm can}/\!\sim
 =\{[\mathrm{PR}]\}.}
 \label{eq:pr35-global-class-singleton}
\end{equation}
Equivalently, the joint exact closure-defect map has one zero orbit.  Every proposed equation has exactly one final disposition: it is an equivalent presentation of the PR equation on the complete declared domain; it has a strictly positive in-class closure defect; or it violates a defining
admissibility guideline and belongs to a different theory.  Hence there is no second inequivalent PR equation architecture and no discrete or continuous free selector in the declared class.
\end{theorem}

\begin{proof}
The minimal even-quartic constraints have the sole coefficient pair $(a_2,a_4)=(1,3/4)$.  Minimal required-channel rank two fixes $\operatorname{Ran}\Pi=\operatorname{span}\{\phi_1,\phi_3\}$.  Signed
pre-projection incidence has the sole zero $|1+\sigma_8|=0$.  The factored hypercharge ideal with the nonzero charge ledger leaves one orbit modulo normalization and up/down relabeling, and the fixed rank-one neutral mass matrix has a one-dimensional photon kernel.  Completeness, nonredundancy, and
normalized incidence give $(N_{\rm gen},[\Phi_3])=(3,[I_3])$.  Finally, the normalized quadratic-response certificate $\rho^2(k-2)^2$ has the sole zero $k=2$.

These unique normal forms are fibers of the acyclic public dependency graph, not independent Cartesian choices.  Every remaining PR-I--PR-III sector map is a fixed inherited relation of the class and therefore contributes no selector coordinate.  Successive composition of fixed maps with singleton fibers remains a singleton.  Operator-equivalent rewritings are removed by $\sim$; an inequivalent in-class proposal has a nonzero component of the joint defect; and a proposal with new operator data, scale, coefficient, order, or field inventory violates the class definition.  These cases exhaust every proposal and prove Eq.~\eqref{eq:pr35-global-class-singleton}.
\end{proof}

The theorem concerns uniqueness of the PR equation architecture.  Different states governed by one fixed field equation are solutions of the same unique theory, not alternative equation architectures. No claim over theories built from different postulates, fields, operator inventories, or power-counting rules is needed for the PR closure statement.
